\documentclass[11pt]{article}
\usepackage[margin=1in]{geometry}
\usepackage{amsmath,amssymb,amsthm}
\usepackage{graphicx}
\usepackage{xcolor}
\usepackage{microtype}
\usepackage{hyperref}
\usepackage{enumitem}

\hypersetup{colorlinks=true,linkcolor=blue!55!black,citecolor=blue!55!black,urlcolor=blue!55!black}
\newtheorem{theorem}{Theorem}
\newtheorem{lemma}{Lemma}
\newcommand{\conv}{\operatorname{conv}}
\newcommand{\vol}{\operatorname{vol}}
\newcommand{\Sph}{S^{n-1}}

\title{Classification of the Mixed Radial--Support\\Composition Inequality as Printed}
\author{Candidate full solution to the second alternative in arXiv:2005.11253, Problem 4.1}
\date{19 July 2026}

\begin{document}
\maketitle

\begin{center}
\fcolorbox{blue!55!black}{blue!4}{\parbox{0.91\textwidth}{
\textbf{Status.} Candidate full solution of the \emph{second alternative as
printed}.  In dimensions $n\ge2$, it holds for all $K_1,K_2$ exactly for
centered Euclidean balls $T=sB_2^n$.  The printed final term is
$T\circ K_2$, not $T\odot K_2$; this distinction is essential.}}
\end{center}

\section{Source question and scope}

Milman and Rotem define two compositions of convex bodies containing the
origin \cite{milman-rotem}.  If $r_K$ and $h_T$ denote radial and support
functions, then
\begin{align*}
 T\odot K
 &=\conv\{r_T(u)r_K(u)u:u\in\Sph\},\\
 T\circ K
 &=\conv\{h_T(u)r_K(u)u:u\in\Sph\}.
\end{align*}
Problem 4.1 on source PDF page 13 asks, among other things, for which $T$
the following inequality holds for all $K_1,K_2$:
\begin{equation}\label{eq:printed}
 |T\odot(K_1+K_2)|^{1/n}
 \ge |T\odot K_1|^{1/n}+|T\circ K_2|^{1/n}.
\end{equation}
Here and below $|\cdot|$ is $n$-dimensional volume.  The mixture of $\odot$
and $\circ$ on the right is exactly what appears in the source.

\begin{figure}[ht]
  \centering
  \includegraphics[width=\textwidth]{figures/open_problem_crop.png}
  \caption{Problem 4.1 on source PDF page 13.  The last right-hand term in the
  second display is $T\circ K_2$.}
\end{figure}

We use the source convention: a convex body is compact and full-dimensional,
and it is only required to contain the origin.  A separate one-dimensional
statement is included.

\section{Proof intuition}

The asymmetry makes the problem rigid.  Insert $K_1=tB_2^n$ and
$K_2=B_2^n$.  Since $T\odot B_2^n=T$, homogeneity forces $T$ to have at least
the volume of $T\circ B_2^n$.  But $r_T\le h_T$ pointwise, so
$T\subseteq T\circ B_2^n$.  The nested bodies must therefore be equal.
Every radial generator $h_T(u)u$ of the larger body now lies back in $T$,
forcing the reverse inequality $h_T\le r_T$.

Thus $r_T=h_T$.  In dimension at least two this common function is constant:
subdivide a spherical geodesic into very short arcs and compare successive
radii using the support inequality.  Hence $T$ is a centered ball.  The
converse is the ordinary Brunn--Minkowski inequality.

\section{Formal result}

For fixed $T$, write
\[
 R_T(K)=T\odot K,\qquad H_T(K)=T\circ K,
\]
and
\[
 F_R(K)=|R_T(K)|^{1/n},\qquad F_H(K)=|H_T(K)|^{1/n}.
\]

\begin{theorem}\label{thm:classification}
Let $n\ge2$ and let $T\subset\mathbb R^n$ be a convex body containing $0$.
Then \eqref{eq:printed} holds for every pair of convex bodies $K_1,K_2$
containing $0$ if and only if
\[
 T=sB_2^n
\]
for some $s>0$.

For $n=1$, every nondegenerate interval containing the origin satisfies the
printed inequality, always with equality.
\end{theorem}

\begin{proof}[Proof of Theorem \ref{thm:classification}]
Assume first that \eqref{eq:printed} holds for every $K_1,K_2$.  Both
compositions are positively homogeneous in their second argument.  Substitute
$K_1=tB_2^n$ and $K_2=B_2^n$, where $t>0$.  Since
$tB_2^n+B_2^n=(t+1)B_2^n$, the printed inequality becomes
\[
 (t+1)F_R(B_2^n)\ge tF_R(B_2^n)+F_H(B_2^n).
\]
Therefore
\begin{equation}\label{eq:reverse-volume}
 |T|^{1/n}=F_R(B_2^n)\ge F_H(B_2^n),
\end{equation}
where we used $R_T(B_2^n)=T$.

For each $u\in\Sph$, the radial boundary point $r_T(u)u$ belongs to $T$.
Consequently
\[
 r_T(u)=\langle r_T(u)u,u\rangle\le h_T(u).
\]
The defining radial generators therefore give
\begin{equation}\label{eq:inclusion}
 T=R_T(B_2^n)\subseteq H_T(B_2^n).
\end{equation}
Together with \eqref{eq:reverse-volume}, this gives equal volumes.  Both sides
of \eqref{eq:inclusion} are full-dimensional convex bodies, so
\begin{equation}\label{eq:body-equality}
 T=H_T(B_2^n)=\conv\{h_T(u)u:u\in\Sph\}.
\end{equation}
In particular, every generator $h_T(u)u$ in \eqref{eq:body-equality} belongs
to $T$.  By maximality of the radial function, $h_T(u)\le r_T(u)$.  Combined
with the reverse pointwise inequality already proved, this yields
\begin{equation}\label{eq:radial-support}
 r_T(u)=h_T(u)=:\rho(u)\qquad(u\in\Sph).
\end{equation}

If $\rho(u)=0$, then applying the next inequality in both directions shows
that $\rho$ vanishes in a spherical neighborhood of $u$.  Because $\rho=h_T$
is continuous, its zero set is also closed.  The sphere is connected for
$n\ge2$, and $T$ is full-dimensional, so $\rho$ has no zeros.

It remains to show that $\rho$ is constant.  Since
$\rho(u)u\in T$, the support inequality in direction $v$ gives
\begin{equation}\label{eq:neighbor}
 \rho(u)\langle u,v\rangle\le h_T(v)=\rho(v).
\end{equation}
Let $u,v\in\Sph$, and subdivide a shortest spherical geodesic from $u$ to
$v$ into $N$ equal arcs of length $\delta<\pi/2$.  Applying
\eqref{eq:neighbor} in both directions to consecutive subdivision points and
multiplying yields
\[
 (\cos\delta)^N
 \le \frac{\rho(v)}{\rho(u)}
 \le (\cos\delta)^{-N}.
\]
Here $\delta=\alpha/N$, where $\alpha\le\pi$ is the geodesic length.  Since
$(\cos(\alpha/N))^N\to1$, letting $N\to\infty$ gives
$\rho(v)=\rho(u)$.  Thus $\rho\equiv s$ and $T=sB_2^n$.

Conversely, if $T=sB_2^n$, then $r_T=h_T=s$, and therefore
\[
 R_T(K)=H_T(K)=sK.
\]
The desired inequality is exactly the classical Brunn--Minkowski inequality
\cite{gardner}, multiplied by $s$.

Finally let $n=1$.  Write $T=[-a,b]$ and $K=[-c,d]$, with positive endpoint
lengths.  On the two directions of $S^0$, radial and support functions agree,
and direct substitution gives
\[
 R_T(K)=H_T(K)=[-ac,bd].
\]
Its one-dimensional volume is $ac+bd$, which is additive under Minkowski sums
of intervals.  Hence the printed inequality is equality for every interval
$T$ containing the origin.
\end{proof}

\section{Verification and limitations}

The proof is non-computational.  The deterministic script in
\texttt{code/numerical\_reconnaissance.py} checks $22{,}324$ polygon pairs
across support and radial modes and finds failures for several nonball
multipliers.  It was used only during discovery and is not evidence for the
theorem.  It may be run from the packet directory with
\texttt{conda run --no-capture-output -n sandbox python
code/numerical\_reconnaissance.py}.

The theorem resolves the second alternative exactly as typeset in Problem
4.1.  It does \emph{not} address the first, all-$\circ$ inequality.  It also
does not address the symmetric radial variant obtained by replacing the final
$T\circ K_2$ in \eqref{eq:printed} with $T\odot K_2$.  If that replacement was
the authors' intent, the present result should be read as a classification of
the literal statement and as evidence that the typographical distinction is
mathematically consequential.

Run indexes were searched for arXiv:2005.11253, the paper title, composition
of convex bodies, and the mixed radial/support formula.  A bounded arXiv search
on 19 July 2026 found the source but no later paper resolving this exact
as-printed alternative.  Human review should focus on the source-formula scope
and the implication
$T=\conv\{h_T(u)u:u\in\Sph\}\Rightarrow r_T=h_T$.

\begin{thebibliography}{9}

\bibitem{milman-rotem}
V. Milman and L. Rotem,
\emph{Novel view on classical convexity theory},
arXiv:2005.11253v1, 2020.

\bibitem{gardner}
R. J. Gardner,
\emph{The Brunn--Minkowski inequality},
Bull. Amer. Math. Soc. (N.S.) \textbf{39} (2002), 355--405.

\end{thebibliography}

\end{document}
